% Independent supplementary text for C4.03 and C4.06, 2026-11-19.

% Add immediately before the C4.03 double-complex construction.
The cover here is an actual cover in the sheaf-theoretic sense; no homotopy invariance for arbitrary coefficient sheaves is being used. In \(X'=X_{I'}^*\), the upper band \(L\) gives the open subspace \(X_L^*\). Choose, cofinally, every neighborhood \(U_\ell\) for \(\ell\in L\) inside \(X_L^*\). Their union is exactly \(X_L^*\), since they contain every stratum with label in \(L\). Therefore
\[
 \bigcup_{j\in J,\,\ell\in L}(U_j\cap U_\ell)
 =\left(\bigcup_{j\in J}U_j\right)\cap X_L^*
\]
stagewise. Every admissible flag has maximum in \(L\), so its intersection is an actual open subset of this union. This is the exact cover asserted on 150/083, and is independent of the pro-homotopy comparison (69).

% C4.03 proof refinement: include near the local exactness argument.
For exactness one works with the augmented complex of the sheaves obtained by direct image from these open intersections, then restricts to a covering open set and contracts by inserting its label (equivalently by union with that flag). A nonflag union has empty intersection and contributes zero. This local contraction avoids an invalid simplification of the stalk of an open direct image at a point outside its open set. With a flabby resolution, all these open direct images are flabby. Hence the horizontal augmented complexes are acyclic after global sections, and the stated double complex computes the cohomology of the union.

% Add to the C4.06 adapted-tubular construction and comparison.
Normalize each adapted tubular germ to have identity derivative on the intrinsic normal bundle. This is possible by composing with the inverse of its fibrewise normal derivative, and it preserves the normal-crossings arrangement and the flag subbundles. The normalization is essential to a choice-independent comparison: arbitrary normal-bundle automorphisms are not being declared homotopic to the identity. In a local branch decomposition every subbundle in the flag is a coordinate subbundle; the tuple of defining sections therefore respects the whole flag simultaneously. If a kernel block is presented without a chosen line decomposition, retain its coordinates from an ordinary tubular chart and use the defining sections in the quotient block; the inverse-function argument is unchanged.

One also explicitly compares the small-sector germ with the entire punctured bundle. Choose a bundle norm and any positive continuous radius function \(\eta\) on the base. The inclusion
\[
 \{v\in E^*:\|v\|<\eta(\pi v)\}\hookrightarrow E^*
\]
is a homotopy equivalence: the map
\[
 v\longmapsto
 \frac{\eta(\pi v)}{\eta(\pi v)+\|v\|}\,v
\]
lands in that open disk sector, and linear interpolation of its strictly positive scalar factor with \(1\) supplies both inverse homotopies. All coordinate nonvanishing conditions are preserved. Variable-radius disk neighborhoods are cofinal by the tubular chart and paracompactness of the base, so their natural inclusions give an equivalence of the sector prospace with the constant space \(E^*\). Only after this comparison is the canonical quotient \(q:E^*\to Q^*\) applied.
